\section{Introduction}
\label{p:intro}

The (3+1)-dimensional Dirac equation, the Clifford algebra it sits on, and
the spinorial Lorentz generators that organize it, all sit comfortably in
the continuum: they assume a smooth manifold and a continuous symmetry
group. A persistent question, going back to Wilson's lattice fermions and
revisited recently in quantum-walk and quantum-cellular-automata
frameworks~\cite{Arrighi2018,Farrelly2020,FarrellyShort2014}, is whether
this whole relativistic edifice can be obtained as an effective infrared
sector of a discrete underlying dynamics. The answer is, by now, clearly
``yes'' at the level of low-energy effective Hamiltonians, but the
\emph{verification} side has remained a patchwork of analytic
calculations.

% Figure 1: dimensional ladder
\begin{figure}[t]
\centering
\begin{tikzpicture}[
  every node/.style={align=center, font=\small},
  box/.style={draw, rounded corners, minimum width=2.8cm, minimum height=1.7cm,
              align=center, font=\small},
  arr/.style={->, >=Latex, thick, shorten >=2pt, shorten <=2pt},
]
  % Top row: discrete models
  \node[box, fill=blue!8]   (ssh)  at (0,2)
    {SSH 1+1\\[2pt]\(H{=}(t_1{+}t_2\cos k)\sigma_x\)\\\(+ t_2\sin k\,\sigma_y\)};
  \node[box, fill=green!8]  (gr)   at (4.2,2)
    {Graphene 2+1\\[2pt]Tight-binding\\\(\Phi(\vec k)\;\sigma_+\) {+} h.c.};
  \node[box, fill=orange!10] (qw)  at (8.4,2)
    {Honeycomb QW\\[2pt]\(U=\prod_i S_i\,C(\theta)\)};
  \node[box, fill=red!10]   (wd)   at (12.6,2)
    {Weyl/Dirac 3D\\[2pt]\(H_W{=}v\,\vec\sigma\cdot\vec p\)};

  % Bottom: continuum target
  \node[box, fill=purple!10, minimum width=5cm] (dirac) at (6.3,-0.6)
    {Dirac in 3+1 (effective infrared)\\[2pt]
     \(H_D{=}v\,\tau_z\!\otimes\!\vec\sigma\!\cdot\!\vec p \,{+}\, m\,\tau_x\!\otimes\!\bbI\),
     \quad \(H_D^2{=}(v^2|\vec p|^2{+}m^2)\bbI_4\)};

  % Arrows along the top row
  \draw[arr] (ssh) -- node[above, font=\scriptsize]
              {add 2nd dim} (gr);
  \draw[arr] (gr) -- node[above, font=\scriptsize]
              {explicit step} (qw);
  \draw[arr] (qw) -- node[above, font=\scriptsize]
              {3D + chirality} (wd);

  % Arrows from each top box down to Dirac 3+1
  \draw[arr] (ssh.south) .. controls +(0,-0.8) and +(-1,1) .. (dirac.north west);
  \draw[arr] (gr.south) -- ([xshift=-1.5cm]dirac.north);
  \draw[arr] (qw.south) -- ([xshift=1.5cm]dirac.north);
  \draw[arr] (wd.south) .. controls +(0,-0.8) and +(1,1) .. (dirac.north east);

  % Side label
  \node[font=\footnotesize\itshape, right=0.2cm of wd, xshift=-0.3cm, yshift=-0.4cm] {};
\end{tikzpicture}
\caption{The dimensional ladder. Four discrete models (top row) all flow
under linearization plus continuum limit to the same effective Dirac block
in $3{+}1$ (bottom). The algebraic objects that survive the limit
(Clifford algebra, $\Sigma^{\mu\nu}$, chirality balance) are the same at
each rung.}
\label{p:fig:ladder}
\end{figure}


This paper makes three coordinated contributions.

\paragraph{1.~Ladder.}
We assemble a single coherent ladder, dimension by dimension:
\[
\text{SSH in } 1{+}1 \;\longrightarrow\;
\text{graphene in } 2{+}1 \;\longrightarrow\;
\text{honeycomb QW} \;\longrightarrow\;
\text{Weyl/Dirac in } 3{+}1 ,
\]
with the same notation and the same algebraic objects all the way
through.\footnote{The pedagogical ordering is non-trivial: the standard
graphene derivation already contains everything one needs to do the
honeycomb QW, but the QW makes the discrete-to-continuous transition
\emph{explicit}, which is where the relativistic structure earns its
status as an emergent infrared organization.} Throughout, we keep careful
track of the distinction between (i) microscopic exact symmetries of the
discrete step, (ii) effective infrared symmetries of the linearized
Hamiltonian, and (iii) covariant Lorentz / Clifford structures organizing
the infrared limit.

\paragraph{2.~Formal verification.}
Every algebraic identity that the paper claims as a result is paired with
a machine-checked symbolic test. \texttt{sympy} handles closed-form
identities (Clifford anticommutators, $\Sigma^{\mu\nu}$ commutators,
dispersions, Taylor expansions around Dirac points). \texttt{z3}~\cite{deMoura2008}
handles combinatorial and order-theoretic constraints
(Nielsen--Ninomiya balance, isotropy, structural consistency of small
causal sets, bipartite index bounds). Every chapter and every section
carries a comment of the form
\begin{center}\small\texttt{\% verified-by: validators/\textless module\textgreater.py::\textless test\textgreater}\end{center}
that links the claim to the test that backs it. The build pipeline
(\texttt{scripts/build\_book.ps1}) runs \texttt{pytest} before running
\texttt{latexmk}; the PDF is produced \emph{only} if all $146$ tests pass.

\paragraph{3.~Open problems with first-line evidence.}
Two problems are left open: a global classification of the discrete
quasi-energy spectrum on the Brillouin torus (the \emph{spectral} branch),
and a formulation of the underlying \emph{protospace} that does not
depend on the reciprocal-space description (the \emph{structural}
branch). For both we provide a first body of verified machinery
(Sections~\ref{p:open-problems}), including a finite-graph treatment of
chirality balance that does not invoke a torus.

\paragraph{Scope.}
We do not claim that the ladder is unique, nor that the protospace is
fixed by it; we claim that the ladder is internally consistent, that the
relativistic algebra arises as an effective infrared organization rather
than as a microscopic symmetry, and that the verification methodology is
fine-grained enough to be useful beyond this corpus.
